%%%%%
%
%% Dans ce fichier, les �ventuels accents sont cod�s
%% en ISO-8859-15 (aka latin9)
%% La ligne suivante donne trois exemples 
%% e accent aigu : �, e accent grave : �, e-dans-l'o : �
\section{Caract�ristique int�grale de la turbulence th�orique selon Foken avec altitude variable}

\itemstyle{Op\'erateur~:} \operfmt{ITC\_model\_Foken\_zvar}
\index{ITC\_model\_Foken\_zvar}

%%%%%OBJET%%%%
\itemstyle{Objet~:}
Calcul des caract�ristiques int�grales de la turbulence en utilisant 
le mod�le propos� par Thomas Foken avec une altitude variable.

%%%%%ARGUMENTS%%%%
\bigskip
%% Dans ce fichier, les �ventuels accents sont cod�s
%% en ISO-8859-15 (aka latin9)
%% La ligne suivante donne trois exemples 
%% e accent aigu : �, e accent grave : �, e-dans-l'o : �
%%%%%ENTREES%%%%
\begin{tabarg}
\multicolumn{3}{l}{\titretab {Entr\'ees}}\\
\hline
$U_*$ & Vecteur, \texttt{float} & Vitesse de frottement (m/s)\\
$L$ & Vecteur, \texttt{float} & Longueur de Monin-Obukhof (m)\\
$Z$ & Vecteur, \texttt{float} & Hauteur de mesure (m)\\
$Hc$ & Coeff., \texttt{float} & Hauteur de la canop�e (m)\\
$\lambda$  & Coeff., \texttt{float} & Latitude (degr�s)\\
\end{tabarg}

\medskip
%%%%%SORTIES%%%%
\begin{tabarg}
\multicolumn{3}{l}{\titretab {Sorties}}\\
\hline
$ITC_U$ & Vecteur, \texttt{float} & Caract�ristique int�grale de la turbulence selon U\\
$ITC_W$ & Vecteur, \texttt{float} & Caract�ristique int�grale de la turbulence selon W \\
$ITC_T$ & Vecteur, \texttt{float} & Caract�ristique int�grale de la turbulence selon T \\
\end{tabarg}


%%%%%FORMULE%%%%
\itemstyle{Formule}

%% Formule de l'operateur : ITC_model_Foken_zvar
%% Dans ce fichier, les �ventuels accents sont cod�s
%% en ISO-8859-15 (aka latin9)
%% La ligne suivante donne trois exemples 
%% e accent aigu : �, e accent grave : �, e-dans-l'o : �

$$
ITC = c_1.  \left\vert \frac{Z}{L} \right\vert ^{c_2};
$$

avec $c_1$ et $c_2$ d�pendant de $Z/L$ suivant le tableau ci-dessous : \\

\begin{center}
\begin{tabular}{llll}
\hline 
Param�tre & $Z/L$ & $c_1$ & $c_2$ \\ 
\hline 
$ITC_W$ & $ \vert Z/L \vert < 0,032 $ & 1,3 &  0 \\ 
  & $\vert Z/L \vert \geq 0,032$ & 2,0 & 1/8 \\ 
\hline 
$ITC_U$ & $\vert Z/L \vert < 0,032$ & 2,7 &  0 \\ 
  & $\vert Z/L \vert \geq 0,032$ & 4,15 & 1/8 \\
\hline 
$ITC_T$ & $0,02 \leq Z/L  < 1$ & 1,4 & -1/4 \\ 
 & $0,02 > Z/L > -0,062 $ & 0,5 & -1/2 \\
 & $-0,062 > Z/L > -1 $ & 1,0 & -1/4 \\
 & $-1 > Z/L  $ & 1,0 & -1/3 \\
 & $Z/L \geq 1 $ & 1/4 & -1/3 \\
\hline 
\end{tabular}
\end{center}
\vspace{0cm}
Si $-0,2 < Z/L < 0,4 $ alors :

$$
ITC_W  = 0,21 * log(f / U_*) + 3,1 
$$
$$
ITC_U = 0,44 * log(f / U_*) + 6,3
$$

avec $f = 2.\frac{2\pi}{24x3600}.\sin(\lambda) $


%%%%%AUTEUR%%%%
\itemstyle{Auteur~:}
CNRM/GMEI/TRAMM

%%%%%REFERENCE%%%%
\itemstyle{R\'ef\'erences~:}

Micro-meteorology, T. Foken, 2008 \cite{MICROMETEO-FOKEN}\\
Sources du module QC\_QA de TK3 disponible en ligne : \\
https://www.bayceer.uni-bayreuth.de/mm/de/software/software/software\_dl.php

